\chapter{DishBrain}
\chaptermark{DishBrain}
\label{sec:dishbrain}
\begin{chaptergoals}
\item як DishBrain кодує м'яч місцем і частотою, а ракетку читає як різницю лічильників;
\item чому лічба у вікні $10\ms$ шумна, тож мережа може зсунути лише середнє;
\item як мережа, яку струшують лише після промахів, дрейфує до налаштувань, що рідше помиляються;
\item що виміряли, про що сперечаються і як перебудувати стенд, щоб перевірити механізм.
\end{chaptergoals}

\section{Мережа в чашці грає в теніс}

DishBrain --- так автори назвали стенд, на якому культура живих нейронів, вирощена на чипі з електродами, грає в спрощений теніс з однією ракеткою~\cite{Kagan2022}. Це найобговорюваніший дослід із машиною з живих нейронів (огляд таких машин --- у розділі~\ref{sec:livingmachine}), і для нашої книги він особливо зручний. Тут маленька мережа доступна вся: кожен вхід задає експериментатор, кожен вихід записується. Усі питання попередніх розділів --- як записано сигнал, хто навчає, що бачить щуп --- можна поставити мережі, яка не має ні очей, ні м'язів, ні \nt{DOPA}-нейронів. Розберімо дослід так, як інженер розбирає чужий стенд: схема, вхід, вихід, учитель, результат, суперечки.

\section{Стенд}

Нейрони беруть із кори зародка миші --- близько $800\,000$ клітин на чип --- або вирощують зі стовбурових клітин людини, близько мільйона на чип. Кілька тижнів вони ростуть на чипі, проростають одна в одну й починають самі видавати імпульси та пачки. На майданчику $8$~мм$^2$ під ними лежить $26\,000$ платинових електродів; одночасно можна записувати до $1024$ із них, а стимулювати на практиці --- через вісім.

Навколо культури замкнено петлю (рис.~\ref{fig:dishloop}). Комп'ютер моделює гру: м'яч летить, відбивається від стін, ракетка рухається вгору і вниз. Положення м'яча перетворюється на струм на восьми «чутливих» електродах. Імпульси з двох «рухових» ділянок підраховують вікнами по $10\ms$, і вони рухають ракетку. Після кожного удару чи промаху культура отримує ще один стимул --- зворотний зв'язок. Вікна по $10\ms$ --- це такт комп'ютера стенда, а не культури: сама мережа, як і раніше, працює асинхронно, як і всі мережі цієї книги.

\begin{figure}[t]
\centering
\begin{tikzpicture}[>=Stealth,
  blk/.style={draw,very thick,rounded corners=3pt,align=center,font=\small,minimum height=1.2cm},
  lab/.style={font=\scriptsize,align=center}]
\node[blk,fill=blue!6,text width=3.4cm] (C) at (0,0) {культура на чипі\\\scriptsize $\sim10^6$ нейронів, $26\,000$ електродів};
\node[blk,fill=orange!10,text width=3.0cm] (G) at (7.2,0) {комп'ютер:\\гра в теніс};
\draw[->,very thick,blue!60!black] (G.north west) to[bend right=25] node[lab,above] {м'яч: \emph{місце} --- який із 8 електродів,\\\emph{частота} --- 4--40 імп/с} (C.north east);
\draw[->,very thick,red!70!black] (C.south east) to[bend right=25] node[lab,below] {ракетка: лічба імпульсів\\ділянок «вгору» і «вниз» за $10\ms$} (G.south west);
\draw[->,thick,densely dashed,green!45!black] (G.east) -- ++(0.9,0) |- ++(0,-2.6) -| node[lab,pos=0.25,below] {відбив: передбачуваний стимул; промах: непередбачуваний} (C.south);
\end{tikzpicture}
\caption{Петля DishBrain. Синя стрілка --- вхід: положення м'яча записано кодом місця і частотним кодом (розділ~\ref{sec:code}). Червона --- вихід: різниця лічильників двох ділянок рухає ракетку. Зелена штрихова --- зворотний зв'язок після кожного розіграшу, єдиний «учитель» стенда.}
\label{fig:dishloop}
\end{figure}

\section{Вхід і вихід}

\emph{Вхід} записано одразу двома кодами розділу~\ref{sec:code}. Те, який із восьми електродів стимулюється, каже, на якій висоті м'яч: це код місця, як у датчиків розділу~\ref{sec:sensors}. Те, як часто йдуть імпульси стимуляції, каже, наскільки м'яч далеко: $4$ на секунду, коли він біля дальньої стіни, і лінійно більше, до $40$ на секунду, коли він біля стіни ракетки. Це частотний код; амплітуда імпульсів м'яча --- $75$~мВ. Культура не «знає» заздалегідь жодного з цих кодів: код задав експериментатор, і мережа має навчитися його читати.

\emph{Вихід} читають так само, як в інтерфейсах розділу~\ref{sec:debugging}. Імпульси однієї рухової ділянки штовхають ракетку вгору, іншої --- вниз; у найпростішому записі за кожне вікно
\begin{empheq}[box=\keybox]{equation}
\Delta p \;=\; c\,\bigl(n_\uparrow - n_\downarrow\bigr),
\label{eq:dishreadout}
\end{empheq}
де $n_\uparrow$, $n_\downarrow$ --- кількості імпульсів двох ділянок за $10\ms$, а $c$ --- коефіцієнт стенда. Це двопровідний код розділу~\ref{sec:glutcomputer}: знак руху несе не знак сигналу, а те, який із двох проводів гучніший.

Погляньмо на шум такого виходу. Якщо ділянка видає загалом $R$ імпульсів на секунду, за вікно $T=10\ms$ їх у середньому $RT$, і за~\eqref{eq:rateerror} відносний розкид дорівнює $1/\sqrt{RT}$. При $R=1000$ імпульсів на секунду це близько $30$ відсотків у кожному вікні; при $R=100$ --- понад сто. Ракетка неминуче тремтить, і навчитися мережа може лише в один спосіб --- зсунути \emph{середню} різницю лічильників у потрібний бік. Це ті самі закони, що обмежують будь-який інтерфейс розділу~\ref{sec:debugging}, тільки тепер по обидва боки петлі.

\section{Учитель без дофаміну}

Найцікавіше в стенді --- учитель. У культурі нейронів кори немає \nt{DOPA}-нейронів, і жодного сигналу «краще, ніж очікувалося» стенд не надсилає. Зворотний зв'язок інший. Відбила мережа м'яч --- на всі вісім чутливих електродів одразу подається короткий \emph{передбачуваний} стимул: імпульси по $75$~мВ, $100$ на секунду впродовж $0.1\,\text{с}$. Схибила --- чотири секунди триває стимул, який автори називають \emph{непередбачуваним}: імпульси по $150$~мВ, $5$ на секунду, у випадкові моменти на випадково вибрані електроди; далі чотири секунди паузи без стимуляції, і м'яч стартує знову~\cite{Kagan2022}.

Автори виходили з гіпотези, що мережа діє так, щоб її вхід ставав передбачуваним~\cite{Friston2010}. Для інженера зміст простий: культура має рівно один спосіб позбутися чотирьох секунд випадкового шуму --- відбити м'яч. Ту саму ідею ми зустрічали в економії імпульсів розділу~\ref{sec:code} і в передбаченні кадру в розділі~\ref{sec:recognition}.

Але чи потрібна тут саме передбачуваність? Можливий і грубіший механізм. Нехай непередбачуваний стимул після промаху просто \emph{перемішує} недавні зміни в мережі, а спокійний після удару їх не зачіпає. Опишімо налаштування мережі одним вектором $\boldsymbol\psi$, і нехай у налаштуванні $\boldsymbol\psi$ мережа хибить з імовірністю $P_{\text{пр}}(\boldsymbol\psi)$. Якщо поштовхи симетричні --- з $\boldsymbol\psi$ у $\boldsymbol\psi'$ так само ймовірно, як назад, --- то в усталеному режимі потік виходів із кожного налаштування врівноважений потоком входів, і частка часу, яку мережа проводить у налаштуванні $\boldsymbol\psi$, дорівнює
\begin{empheq}[box=\keybox]{equation}
\rho(\boldsymbol\psi) \;\propto\; \frac{1}{P_{\text{пр}}(\boldsymbol\psi)} .
\label{eq:dishdensity}
\end{empheq}
Налаштування, яке хибить удесятеро рідше, утримується вдесятеро довше. Ніхто не повідомляє мережі, в який бік змінювати ваги: добрі налаштування просто рідше руйнуються.

Ми перевірили це на моделі, звівши культуру до двох чисел: ракетка приходить у точку $a\,y+b$, коли м'яч приходить на висоту $y$, і відбиває, якщо схибила менш ніж на $0.1$. Після промаху обидва числа отримують випадковий поштовх, після удару лишаються незмінними. Частка відбитих м'ячів за $2000$ розіграшів зросла з $0.21$ у перших двохстах до $0.38$ в останніх п'ятистах (сорок «культур», рис.~\ref{fig:dishmodel}). Якщо поштовх давати після \emph{кожного} розіграшу, у зворотному зв'язку немає відомостей, і частка лишається близько $0.12$; якщо поштовхів немає зовсім, вона стоїть на $0.19$. Дослід на стенді з цим узгоджується: значуще зростання всередині сеансу було тільки за повного зворотного зв'язку. Коли замість стимулів була тиша, культура наприкінці сеансу все ж грала краще, ніж без зворотного зв'язку, але з меншим ефектом~\cite{Kagan2022}; зауважмо, що і в цій умові після промаху йде довга пауза, а після удару --- коротка, тож різниця між наслідками не зникає цілком.

\begin{figure}[t]
\centering
\begin{tikzpicture}[x=1cm,y=5cm]
\draw[->,gray!70!black] (0,0) -- (10.6,0);
\draw[->,gray!70!black] (0,0) -- (0,0.5) node[above,font=\small,black] {частка відбитих};
\foreach \v in {0.1,0.2,0.3,0.4}{\draw[gray!60!black] (-0.06,\v) -- (0.06,\v); \node[left,font=\scriptsize] at (0,\v) {\v};}
\foreach \x/\f/\l/\lab in {1.8/0.21/0.38/{шум після\\промаху},5.2/0.13/0.11/{шум після\\кожного розіграшу},8.6/0.19/0.19/{без шуму}}{
  \fill[blue!35] (\x-0.8,0) rectangle (\x-0.05,\f);
  \fill[red!60!black] (\x+0.05,0) rectangle (\x+0.8,\l);
  \node[below,font=\scriptsize,align=center] at (\x,-0.01) {\lab};
  \node[above,font=\scriptsize] at (\x-0.425,\f) {\f};
  \node[above,font=\scriptsize] at (\x+0.425,\l) {\l};}
\fill[blue!35] (7.4,0.47) rectangle (7.7,0.495); \node[right,font=\scriptsize] at (7.75,0.482) {перші 200};
\fill[red!60!black] (7.4,0.41) rectangle (7.7,0.435); \node[right,font=\scriptsize] at (7.75,0.422) {останні 500};
\end{tikzpicture}
\caption{Модель «перемішування при промаху» (\texttt{dishbrain\_model.py}): частка відбитих м'ячів на початку і наприкінці $2000$ розіграшів, середнє за сорока моделями-культурами. Навчання йде, тільки якщо поштовх настає після промаху, але не після удару, --- як і в досліді, де значуще зростання за сеанс було лише за повного зворотного зв'язку.}
\label{fig:dishmodel}
\end{figure}

\begin{keyblock}
\begin{proposition}[Учитель-перемішувач]
\label{prop:dishscramble}
Якщо невдача струшує мережу, а вдача залишає її в спокої, мережа сама дрейфує до налаштувань, які рідше помиляються: час, проведений у налаштуванні, обернено пропорційний імовірності невдачі~\eqref{eq:dishdensity}. Для такого навчання не потрібні ні сигнал винагороди, ні його знак, ні передбачуваність як така --- досить, щоб зворотний зв'язок після вдачі та після невдачі різнився. Зміну поведінки культури виміряно; який механізм працює в чашці --- передбачення входу чи перемішування, --- не встановлено.
\end{proposition}
\end{keyblock}

Порівняйте з розділом~\ref{sec:dofa}. Там шум теж слугував пошуком, але дофамін мав знак: помилка $\delta$ казала, в який бік зсунути ваги, і крок ішов за градієнтом~\eqref{eq:perturbation}. Тут знака немає --- є лише «залишити» або «струснути». Такий пошук грубіший і повільніший, але вимагає від мережі менше: ні \nt{DOPA}-нейронів, ні критика, ні сліду, розтягнутого на секунди.

\section{Що виміряли і про що сперечаються}

Кожна гра тривала $20$ хвилин; порівнювали перші $5$ хвилин з останніми $15$. У культур із повним зворотним зв'язком серії відбитих м'ячів стали довшими, а миттєвих промахів поменшало; у контролях без клітин, без гри та з комп'ютерною «ракеткою» нічого подібного не було. Культури з клітин людини в середньому відбивали довше за мишачі. Покращення статистично надійне --- на дев'яти мишачих і одинадцяти людських культурах, у сотні з гаком ігор кожної групи, --- але невелике: мережа не стала добрим гравцем, вона стала трохи кращою за випадкову. Надійне воно всередині сеансу; стійкого навчання від сеансу до сеансу впродовж кількох днів автори не отримали~\cite{Kagan2022}. У наступній роботі тієї самої групи виявили, що під час гри активність культури наближається до \emph{критичного} режиму --- межі між згасанням і лавиною, того самого життя на межі, що в розділі~\ref{sec:gaba}, --- і що ближче, то краща гра~\cite{Habibollahi2023}.

Головну суперечку викликали не цифри, а слова. Заголовок статті стверджував, що нейрони в чашці «виявляють розумність» (sentience), і група дослідників заперечила: зміна поведінки в замкненій петлі --- ще не розум і не відчуття, а висновки треба формулювати скромніше~\cite{Balci2023}. З інженерного погляду відкриті два питання, і обидва --- про механізм. Перше: чи навчається мережа, тобто чи змінюються її ваги надовго, чи зворотний зв'язок лише зсуває її поточний стан? Друге: чи потрібна саме передбачуваність, чи вистачає будь-якої різниці між відгуком на вдачу і на невдачу, як у твердженні~\ref{prop:dishscramble}? Стенд, на якому доступний кожен електрод, дає змогу відповісти на обидва --- і це, мабуть, його головна цінність.

\section{Специфікація стенда: як його відтворити}
\label{sec:dishspec}

Нижче зібрано все, що потрібно, щоб зібрати такий стенд наново: обладнання, петля, кодування входу, читання виходу, зворотний зв'язок і протокол досліду. Кожен параметр позначено джерелом. «Стаття» --- значення взято з тексту роботи~\cite{Kagan2022}. «Вибір» --- у статті його немає, і для відтворення його доведеться задати самому; ми пропонуємо типове значення і кажемо, як його перевірити.

\subsection*{Обладнання і культура}

\begin{center}
\footnotesize
\setlength{\tabcolsep}{4pt}
\renewcommand{\arraystretch}{1.15}
\begin{tabular}{>{\raggedright}p{4.0cm}>{\raggedright}p{6.9cm}>{\raggedright\arraybackslash}p{1.6cm}}
\toprule
Параметр & Значення & Джерело \\
\midrule
чип & матриця MaxOne: $26\,000$ платинових електродів на майданчику $8$~мм$^2$ & стаття \\
запис & до $1024$ каналів одночасно, $20\,000$ відліків на секунду & стаття \\
стимуляція & до $32$ електродів технічно, у досліді --- $8$ незалежних & стаття \\
клітини & кора зародка миші; нейрони людини зі стовбурових клітин (два способи вирощування); клітини HEK293T (не нейрони) для контролю & стаття \\
висів & близько $10^6$ клітин на чип & стаття \\
вік культури на час досліду & миша: від $\sim10$ діб; людина: $14$--$24$ доби або $\sim80$ діб залежно від способу вирощування & стаття \\
\bottomrule
\end{tabular}
\end{center}

\subsection*{Петля і час}

Петля працює вікнами по $10\ms$ ($200$ відліків): за вікно підраховують імпульси на електродах рухових ділянок, гра оновлюється, і видається чергова стимуляція. Затримка від імпульсу в культурі до відповідного стимулу --- близько $5\ms$. Під час гри сигнал кожного електрода проходить фільтр верхніх частот Бесселя 2-го порядку з частотою зрізу $100$~Гц; модуль сигналу згладжується фільтром нижніх частот Бесселя 1-го порядку з частотою $1$~Гц, і поріг імпульсу пропорційний цьому згладженому модулю. Поріг у шість стандартних відхилень фону стаття вказує для окремих вимірювань активності, а не для гри. Щоб стимуляцію не зараховували як відповідь культури, усі сигнали гасять, коли одночасно надходить понад $15$ великих, вищих за $75$~мВ, імпульсів. Усе це --- стаття.

\subsection*{Кодування м'яча}

\begin{center}
\footnotesize
\setlength{\tabcolsep}{4pt}
\renewcommand{\arraystretch}{1.15}
\begin{tabular}{>{\raggedright}p{3.4cm}>{\raggedright}p{7.5cm}>{\raggedright\arraybackslash}p{1.6cm}}
\toprule
Параметр & Значення & Джерело \\
\midrule
чутлива ділянка & $8$ електродів, розташованих так, щоб сусідні електроди означали сусідні положення м'яча & стаття \\
код місця & номер електрода $k=1,\dots,8$ задає положення м'яча відносно центру ракетки & стаття \\
яка координата & вертикальне положення м'яча; для відтворення --- відносно ракетки, $k=\lceil 8(y-p+\tfrac12)\rceil$ при $y-p\in[-\tfrac12,\tfrac12]$ & стаття; формула --- вибір \\
частотний код & частота стимуляції від $4$ до $40$~імп/с несе відстань до ракетки & стаття \\
напрям шкали & $4$~імп/с, коли м'яч біля протилежної стіни, і лінійно до $40$~імп/с, коли він біля стіни ракетки: $f=4+36\,(1-x/L)$, $x$ --- відстань до стіни ракетки, $L$ --- ширина корту & стаття \\
форма імпульсу & короткий прямокутний двофазний: спершу додатна, потім від'ємна фаза & стаття \\
амплітуда & $75$~мВ; обрано так, щоб не змушувати спрацьовувати загальмовані клітини & стаття \\
тривалість фаз & не вказано; взяти стандартне налаштування системи стимуляції і не змінювати його впродовж досліду & вибір \\
\bottomrule
\end{tabular}
\end{center}

\subsection*{Читання ракетки}

У статті дві рухові ділянки: імпульси в першій рухають ракетку вгору, у другій --- вниз; внесок ділянок зважують, щоб урахувати різну щільність клітин і різну активність~\cite{Kagan2022}. Точної формули не наведено. Для відтворення беремо правило~\eqref{eq:dishreadout}, $\Delta p=c\,(n_\uparrow-n_\downarrow)$ за кожне вікно $10\ms$, з вагою, що вирівнює ділянки за їхньою середньою активністю в спокої (вибір): $n_\uparrow$ і $n_\downarrow$ ділять на середню лічбу своєї ділянки за вікно, виміряну перед грою. Коефіцієнт $c$ обираємо так, щоб різниця в одну середню лічбу зсувала ракетку на $1\%$ висоти корту за вікно (вибір); тоді стала перевага однієї ділянки проводить ракетку через увесь корт за $1$~с.

Розташування ділянок на чипі в тексті статті координатами не задано; є лише схема. Для відтворення досить трьох груп електродів, що не перетинаються: вісім чутливих у центрі та по групі записувальних електродів з обох боків, одна --- «вгору», друга --- «вниз» (вибір).

\subsection*{Гра}

Розмірів корту, довжини ракетки і швидкості м'яча в статті не вказано; сказано лише, що м'яч летить зі сталою швидкістю і відбивається від стін. Для відтворення (вибір): корт $1\times1$, ракетка на правій стінці завдовжки $0.2$ (влучання при $|y-p|<0.1$, як у моделі~\texttt{models/dishbrain\_model.py}), м'яч перетинає корт за $1$~с. Промах без жодного відбитого м'яча в розіграші вважають «ейсом»; розіграш, довший за три удари поспіль, --- «довгим» (стаття).

\subsection*{Зворотний зв'язок}

\begin{center}
\footnotesize
\setlength{\tabcolsep}{4pt}
\renewcommand{\arraystretch}{1.15}
\begin{tabular}{>{\raggedright}p{2.6cm}>{\raggedright}p{8.3cm}>{\raggedright\arraybackslash}p{1.6cm}}
\toprule
Подія & Що подається & Джерело \\
\midrule
удар & передбачуваний стимул: усі $8$ чутливих електродів одночасно, $75$~мВ, $100$~імп/с, $100\ms$; звичайна стимуляція м'яча на цей час переривається & стаття \\
промах & непередбачуваний стимул: $150$~мВ, $5$~імп/с, $4$~с, у випадкові моменти на випадково вибрані з $8$ електроди; потім $4$~с без стимуляції, і м'яч стартує у випадковому напрямку & стаття \\
закон випадковості & не вказано; для відтворення --- моменти імпульсів як пуассонівський потік із середньою частотою $5$~імп/с & вибір \\
\bottomrule
\end{tabular}
\end{center}

\subsection*{Протокол і контроль}

Один сеанс триває $20$~хв; порівнюють перші $5$ і останні $15$~хв (стаття). Три міри результату: середня довжина розіграшу, частка ейсів і частка довгих розіграшів. Умови досліду (стаття):
\begin{itemize}
\item \emph{стимул} --- повна петля, як описано вище;
\item \emph{тиша} --- замість стимулу удару чи промаху стільки ж часу без жодної стимуляції, потім м'яч стартує у випадковому напрямку;
\item \emph{без зворотного зв'язку} --- м'яч після промаху просто відбивається й летить далі, стимуляція положення м'яча триває;
\item \emph{спокій} --- гра йде і ракетка рухається за активністю, але стимуляції немає зовсім;
\item \emph{контролі} --- чип лише із середовищем; клітини HEK293T; «культура в комп'ютері», симуляція замість клітин.
\end{itemize}
Обсяг вибірки в основній серії: $9$ культур миші ($101$ сеанс), $11$ культур людини ($138$), $6$ чипів із середовищем ($80$), $20$ культур у спокої ($42$), $3$ симуляції ($38$), разом $399$ сеансів. У серії про зворотний зв'язок --- $486$ сеансів за $3$ дні по $3$ сеанси на день, умови «стимул», «тиша» і «без зворотного зв'язку» ($n=27$, $17$ і $15$). Зростання середньої довжини розіграшу від перших $5$ хвилин до останніх $15$ отримано лише в живих культур. У серії про зворотний зв'язок значуще зростання в часі --- тільки в умові «стимул»; наприкінці сеансу умова «тиша» все ж перевершувала «без зворотного зв'язку», але з меншим ефектом, а стійкого навчання від сеансу до сеансу за три дні не було~\cite{Kagan2022}.

\subsection*{Цикл стенда крок за кроком}

Кожні $10\ms$ комп'ютер стенда робить таке.
\begin{enumerate}
\item Зчитує кількість імпульсів $n_\uparrow$, $n_\downarrow$ у двох рухових ділянках за минуле вікно і нормує їх на середню лічбу ділянки в спокої.
\item Зсуває ракетку на $\Delta p=c\,(n_\uparrow-n_\downarrow)$ і обмежує її краями корту.
\item Зсуває м'яч на один крок; відбиває його від верхньої, нижньої та лівої стінок.
\item Якщо м'яч дійшов до правої стінки: при $|y-p|<0.1$ --- удар, лічба розіграшу зростає на один, подається стимул удару ($100\ms$); інакше --- промах, розіграш записується, подається стимул промаху ($4$~с), потім $4$~с паузи, і м'яч стартує знову.
\item Інакше обирає електрод $k$ за вертикальним положенням м'яча і частоту $f$ за відстанню до стіни ракетки та подає на електрод $k$ двофазні імпульси $75$~мВ із частотою $f$ до наступного вікна.
\end{enumerate}
Цього досить, щоб зібрати стенд, сумісний з опублікованим, і перевірити головне питання розділу: чи навчає культуру передбачуваність, чи проста різниця між відповіддю на удар і на промах. Для цього до трьох умов статті варто додати четверту, якої в ній немає: після промаху --- \emph{передбачуваний} стимул тієї самої тривалості й енергії, що й непередбачуваний. Якщо культура навчається і так, річ у різниці, а не в передбачуваності.


\section{Підсумок}

\begin{table}[t]
\centering
\footnotesize
\setlength{\tabcolsep}{4pt}
\renewcommand{\arraystretch}{1.15}
\begin{tabular}{>{\raggedright}p{2.1cm}>{\raggedright}p{4.2cm}>{\raggedright}p{1.9cm}>{\raggedright\arraybackslash}p{4.6cm}}
\toprule
Вузол стенда & Як улаштований & Розділ книги & Що відомо \\
\midrule
вхід & місце (який із 8 електродів) і частота (4--40 імп/с) & \ref{sec:code}, \ref{sec:sensors} & задано експериментатором \\
вихід & різниця лічильників двох ділянок за $10\ms$~\eqref{eq:dishreadout} & \ref{sec:glutcomputer}, \ref{sec:debugging} & задано експериментатором; шум за~\eqref{eq:rateerror} \\
учитель & передбачуваний стимул після удару, непередбачуваний після промаху; дофаміну немає & \ref{sec:dofa} & значуще зростання за сеанс --- тільки за повного зворотного зв'язку; за тиші ефект менший; від сеансу до сеансу нестійкий \\
механізм & передбачення входу або перемішування при промаху~\eqref{eq:dishdensity} & \ref{sec:dofa}, \ref{sec:recognition} & не встановлено; обидва пояснюють результат \\
режим мережі & наближення до критичного режиму під час гри & \ref{sec:gaba} & зв'язок з якістю гри виміряно \\
«розумність» & твердження авторів & --- & не показано; оскаржено \\
\bottomrule
\end{tabular}
\caption{DishBrain очима інженера.}
\label{tab:dishbrain}
\end{table}

DishBrain --- машина з живих нейронів у найпростішому вигляді: вхід задано кодом, вихід прочитано лічбою, учитель --- різниця між відгуком на вдачу і на невдачу (таблиця~\ref{tab:dishbrain}). Мережа без очей, м'язів і дофаміну трішки навчилася грати, і вже це саме по собі робить її випробувальним стендом для ідей книги. Яким би не виявився механізм --- передбачення, перемішування чи щось третє, --- відповідь знайдуть так, як радить розділ~\ref{sec:debugging}: щупом, тестовим сигналом і вимиканням частин на машині, яку нарешті видно цілком.

\section{Задачі}

\begin{exercise}[\lvlA]
\label{ex:22:1}
\emph{Наскільки треба зсунути лічбу.} Дві ділянки видають разом $R_\uparrow+R_\downarrow=2000$ імпульсів на секунду, потоки пуассонівські. (а)~Який відносний розкид лічби однієї ділянки за $10\ms$ при $R=1000$ і $R=100$? (б)~За $1$~с польоту м'яча зсуви~\eqref{eq:dishreadout} додаються. За якої найменшої $R_\uparrow-R_\downarrow$ середній зсув удвічі більший за розкид?
\end{exercise}

\begin{exercise}[\lvlA]
\label{ex:22:2}
\emph{Скільки розіграшів треба, щоб побачити навчання.} Розіграші незалежні, частка відбитих біноміальна. Скільки розіграшів потрібно в кожному з двох періодів, щоб зростання частки відбитих перевищило два стандартні відхилення різниці: (а)~з $0.20$ до $0.25$; (б)~з $0.21$ до $0.38$, як у моделі?
\end{exercise}

\begin{exercise}[\lvlB]
\label{ex:22:3}
\emph{Виведення~\eqref{eq:dishdensity}.} Налаштування $\boldsymbol\psi$ пробігає скінченну множину; при промаху (імовірність $P(\boldsymbol\psi)>0$) мережа переходить у $\boldsymbol\psi'$ із симетричною ймовірністю $K(\boldsymbol\psi,\boldsymbol\psi')$, при ударі лишається. (а)~Доведіть, що $\rho\propto1/P$ --- стаціонарний розподіл. (б)~Яка частка промахів за двох налаштувань із $P=0.9$ і $P=0.1$?
\end{exercise}

\begin{exercise}[\lvlB]
\label{ex:22:4}
\emph{Поглинальна область моделі.} Ракетка приходить у $p=\min\bigl(1,\max(0,ay+b)\bigr)$, м'яч --- на висоту $y\sim U(0,1)$, удар --- якщо $|p-y|<h=0.1$. (а)~Доведіть, що при $|b|<h$ і $|a-1+b|<h$ мережа відбиває всі м'ячі, і знайдіть площу цієї області. (б)~Знайдіть чисельно середню частку відбитих при $a\sim U(-1,1)$, $b\sim U(0,1)$.
\end{exercise}

\begin{exercise}[\lvlB]
\label{ex:22:5}
\emph{Моделювання: паркан навколо налаштувань.} У копії \texttt{dishbrain\_model.py} дайте $200$ культурам по $20\,000$ розіграшів. Яка частка відбитих в останніх $500$? А якщо після кожного поштовху обрізати налаштування прямокутником $a\in[-1,2]$, $b\in[-1,1]$? Поясніть різницю за допомогою задачі~\ref{ex:22:3}.
\end{exercise}

\begin{exercise}[\lvlB]
\label{ex:22:6}
\emph{Проєктування вхідного коду.} М'яч відбито, якщо помилка за висотою менша за $h=0.1$; мережа читає вхід за $T=0.25$~с, до частоти $r$ розрізняється близько $\sqrt{rT}$ рівнів, стимуляція не частіше за $40$ імпульсів на секунду. (а)~Чи вистачить коду місця на восьми електродах? (б)~Чи можна передати висоту частотою на одному електроді?
\end{exercise}

\begin{exercise}[\lvlC]
\label{ex:22:7}
\emph{Дослідження: передбачення чи перемішування?} Узагальніть модель на $d$ налаштовуваних чисел ($p=\sum_{i<d}a_iy^i$). Як зростає з $d$ кількість розіграшів до частки відбитих $0.5$ при перемішуванні та за правилом зі знаком помилки~\eqref{eq:perturbation}? Який дослід на стенді розділив би дві гіпотези?
\end{exercise}
